% =====================================================================
%  1bat-ud2-12.tex  —  HORA 12: On es talla l'aixeta?
%  (sense preàmbul: s'inclou des de main.tex)
% =====================================================================
\horatitol{Sessió 12 — On es talla l'aixeta?}%
{Interpretar l'arrel d'un ajut com el llindar d'exclusió: qui hi té dret i qui no.}

\subsection{Idea clau}

A les sessions 10 i 11 hem après a \textbf{trobar} l'arrel d'un polinomi. Avui
posem aquesta eina al servei d'una decisió real: per a un ajut que disminueix
amb els ingressos, l'arrel marca exactament el llindar a partir del qual ja
no es té dret a l'ajut — «on es talla l'aixeta».

\begin{idea}
\textbf{Llegir l'arrel en context: el llindar d'exclusió.}
Si un ajut es modelitza amb un polinomi $B(x)$ que disminueix quan
augmenten els ingressos $x$:
\begin{itemize}[leftmargin=1.4em, itemsep=2pt]
  \item Si $x$ és \textbf{menor} que l'arrel: $B(x)>0$, encara hi ha dret
        a l'ajut (i en sabem l'import exacte).
  \item Si $x$ és \textbf{igual} a l'arrel: $B(x)=0$, és el llindar exacte.
  \item Si $x$ és \textbf{més gran} que l'arrel: la fórmula dona un valor
        \textbf{negatiu}, que en la realitat no té sentit com a import: cal
        interpretar-lo com «cap dret a l'ajut» ($0\eur$), no com un deute.
\end{itemize}
\end{idea}

\subsection{Exemples}

\begin{exemple}[el llindar oficial de l'ajut (Activitat 2)]
Recuperem l'ajut $B(x)=-20x+600$ de la sessió 10, on $x$ són els ingressos
mensuals, en milers d'euros. Calculem el seu valor per a diversos ingressos:
\begin{center}
\begin{tabular}{lcccc}
\toprule
$x$ (milers €) & 10 & 20 & 30 & 40 \\
\midrule
$B(x)$ ($\eur$) & 400 & 200 & 0 & $-200$ \\
\bottomrule
\end{tabular}
\end{center}
Ja sabem, de la sessió 10, que l'arrel és $x=30$. El valor a $x=40$ surt
negatiu ($-200$), però \textbf{aquest valor no té sentit real}: no es pot
«deure» un ajut. Cal interpretar-lo com a $0\eur$: a partir de $x=30$, ja no
hi ha cap dret a l'ajut.
\end{exemple}

\begin{exemple}[el cas límit: «menys que» o «menys o igual que»?]
Imaginem dues persones amb ingressos molt propers al llindar: \textbf{Berta},
amb $x=29$, i \textbf{Carles}, amb $x=31$.
\[
B(29)=-20\per29+600=20, \qquad
B(31)=-20\per31+600=-20\;\text{(interpretem com a $0\eur$)}.
\]
Berta encara té dret a un ajut de $20\eur$; Carles ja no en té cap dret. Però,
què passa exactament a $x=30$? Si la normativa diu «té dret mentre
$B(x)>0$» (inequació \emph{estricta}), a $x=30$ ja \textbf{no} hi ha dret,
perquè $B(30)=0$ no és estrictament positiu. Si diu «té dret mentre
$B(x)\geq0$», a $x=30$ encara hi ha dret, però per un import de $0\eur$.
Aquest matís —connectat amb les inequacions estrictes i no estrictes de la
UD1— és exactament «on es talla l'aixeta».
\end{exemple}

\subsection{Exercicis resolts}

\begin{exresolt}[taula de perfils]
L'entitat gestora de l'ajut $B(x)=-20x+600$ revisa quatre expedients.
Calcula $B(x)$ per a cada persona i indica si té dret a l'ajut.
\begin{center}
\begin{tabular}{lccc}
\toprule
Persona & Ingressos $x$ (milers €) & $B(x)$ ($\eur$) & Dret a l'ajut? \\
\midrule
Marta & 12 & ? & ? \\
Jordi & 28 & ? & ? \\
Laia  & 30 & ? & ? \\
Pere  & 34 & ? & ? \\
\bottomrule
\end{tabular}
\end{center}
\solucio
Substituïm cada valor de $x$ al polinomi:
\begin{align*}
B(12)&=-20\per12+600=-240+600=360,\\
B(28)&=-20\per28+600=-560+600=40,\\
B(30)&=-20\per30+600=-600+600=0,\\
B(34)&=-20\per34+600=-680+600=-80 \;\text{(interpretem com a $0\eur$)}.
\end{align*}
\resposta{Marta ($360\eur$) i Jordi ($40\eur$) tenen dret a l'ajut; Laia és
exactament al llindar (cas límit, $0\eur$); Pere ja no hi té dret (la
fórmula dona $-80\eur$, que s'interpreta com a $0\eur$). L'aixeta es talla a
$x=30$ milers d'euros.}
\end{exresolt}

\begin{exresolt}[sensibilitat prop del llindar]
Núria declara $x=\num{29,5}$ milers d'euros d'ingressos, i el seu veí Octavi
en declara $\num{30,5}$. Calcula l'ajut de cadascun i digues qui hi té dret.
\solucio
\begin{align*}
B(\num{29,5})&=-20\per\num{29,5}+600=-590+600=10,\\
B(\num{30,5})&=-20\per\num{30,5}+600=-610+600=-10
   \;\text{(interpretem com a $0\eur$)}.
\end{align*}
\resposta{Núria té dret a $10\eur$; Octavi no en té cap dret, tot i que la
diferència d'ingressos entre tots dos és de només $1.000\eur$.}
\end{exresolt}

\subsection{Exercicis per fer}

\begin{exercicis}
\begin{exlist}
  \item Calcula $B(x)=-20x+600$ per a $x=5$, $x=15$, $x=25$ i $x=32$, i
        indica en cada cas si la persona té dret a l'ajut.

  \item Un altre ajut similar es calcula amb $D(x)=-15x+600$. Troba el seu
        llindar d'exclusió (l'arrel de $D(x)$).

  \item Completa la taula per a l'ajut $B(x)=-20x+600$ i indica, per a
        cada persona, si té dret a l'ajut:
    \begin{center}
    \begin{tabular}{lccc}
    \toprule
    Persona & Ingressos $x$ (milers €) & $B(x)$ ($\eur$) & Dret a l'ajut? \\
    \midrule
    Anna   & 8  & & \\
    Biel   & 22 & & \\
    Carme  & 29 & & \\
    Dídac  & 31 & & \\
    \bottomrule
    \end{tabular}
    \end{center}

  \item Un servei de menjador social calcula l'ajut amb $M(x)=-25x+750$, on
        $x$ són els ingressos familiars mensuals, en centenars d'euros.
        Troba el llindar d'exclusió i digues si una família amb $x=28$
        encara hi té dret (i per quin import).

  \item \textbf{(Compte amb el signe)} Dues persones declaren uns ingressos
        gairebé idèntics per a l'ajut $B(x)=-20x+600$: Èlia amb
        $x=\num{29,8}$ i Farid amb $x=\num{30,2}$. Calcula $B(x)$ per a
        cadascun i indica qui té dret a l'ajut. Vés amb compte: a prop del
        llindar, petites diferències en els ingressos canvien completament
        el resultat (i el seu signe).

  \item \textbf{(Raonament)} La normativa de l'ajut $B(x)=-20x+600$ diu que
        «hi ha dret a l'ajut sempre que $B(x)$ sigui estrictament
        positiu». Marina té exactament $x=30$ (el llindar). Explica si
        Marina té dret a l'ajut, i com canviaria la resposta si la
        normativa digués «sempre que $B(x)$ sigui positiu \emph{o zero}».
        Redacta, en dues o tres frases, una conclusió justificada sobre on
        es «talla l'aixeta» d'aquest ajut.
\end{exlist}
\end{exercicis}

% ---------------------------------------------------------------------
\fullsolucions{Sessió 12}
\begin{solucions}
\begin{sollist}
  \item $B(5)=500\eur$ (dret); $B(15)=300\eur$ (dret); $B(25)=100\eur$
        (dret); $B(32)=-40\eur$, interpretat com a $0\eur$ (sense dret).
  \item $-15x+600=0 \Rightarrow x=40$: el llindar d'exclusió és
        $x=40$ milers d'euros.
  \item Anna: $B(8)=440\eur$ (dret). Biel: $B(22)=160\eur$ (dret). Carme:
        $B(29)=20\eur$ (dret). Dídac: $B(31)=-20\eur$, interpretat com a
        $0\eur$ (sense dret).
  \item $-25x+750=0 \Rightarrow x=30$ (centenars d'euros, és a dir
        $3.000\eur$). Com que $28<30$, la família encara hi té dret:
        $M(28)=-25\per28+750=50\eur$.
  \item $B(\num{29,8})=-20\per\num{29,8}+600=4\eur$ (Èlia té dret). \;
        $B(\num{30,2})=-20\per\num{30,2}+600=-4\eur$, interpretat com a
        $0\eur$ (Farid no en té dret).
  \item Amb la normativa estricta ($B(x)>0$), Marina \textbf{no} té dret,
        perquè $B(30)=0$ no és estrictament positiu. Amb la normativa no
        estricta ($B(x)\geq0$), Marina \textbf{sí} que hi té dret, encara
        que per un import de $0\eur$. Conclusió oberta (exemple vàlid):
        «L'ajut es talla exactament a $30.000\eur$ d'ingressos; per sota
        s'hi té dret, per sobre no, i el cas exacte depèn de si la norma
        accepta la igualtat o no».
\end{sollist}
\end{solucions}
